% 第二章：导数与微分
\chapter{手稿}
\section{杂题收录}
\begin{example}
    已知$x,y>0$，且$x^2+9y^2=12$，则$\dfrac{x+2}{y+1}-3x$的最小值为\underline{\hspace{2cm}}.
\end{example}
\begin{solution}
记
\[
F=\frac{x+2}{y+1}-3x.
\]
因为 $y>0$，要证 $F\ge -6$，等价于
\[
x+2+(6-3x)(y+1)\ge 0,
\]
即
\[
8-2x+6y-3xy\ge 0.
\]
由 $x^2+9y^2=12$ 得
\[
y=\frac{\sqrt{12-x^2}}{3},\qquad 0<x<2\sqrt3.
\]
代入上式，需证
\[
8-2x+2\sqrt{12-x^2}-x\sqrt{12-x^2}\ge0.
\]
当 $0<x\le2$ 时，左式中 $8-2x>0$ 且 $2\sqrt{12-x^2}-x\sqrt{12-x^2}\ge0$，不等式成立。
当 $x>2$ 时，上式等价于
\[
2(4-x)\ge (x-2)\sqrt{12-x^2}.
\]
两边均非负，平方后化为
\[
4(4-x)^2-(x-2)^2(12-x^2)=(x^2-2x-4)^2\ge0.
\]
等号在
\[
x=1+\sqrt5,\qquad y=\frac{\sqrt5-1}{3}
\]
时取得。因此最小值为
\[
-6.
\]
\end{solution}
\newpage
\begin{example}[（高考题改编）]
    （单选）设代数式$\displaystyle T = \dfrac21\cdot\dfrac{4}{3}\cdot\dfrac{6}{5}\cdots\dfrac{200}{199}$，则（~~~~）

    \begin{tabular}{@{}llll@{}}
        A.$14< T < 16$&B.$16 < T < 18$&C.$18< T < 20$&D.$20< T < 22$
    \end{tabular}
\end{example}
\begin{solution}
    由 Wallis 乘积的标准估计，
    \[
    \frac{1}{2n+1}\left(\frac{(2n)!!}{(2n-1)!!}\right)^2<\frac{\pi}{2}<\frac{1}{2n}\left(\frac{(2n)!!}{(2n-1)!!}\right)^2.
    \]
    取 $n=100$，并记 $T=\dfrac{(200)!!}{199!!}$，则
    \[
    100\pi<T^2<\frac{201\pi}{2}.
    \]
    而 $17^2=289<100\pi<\frac{201\pi}{2}<324=18^2$，所以
    \[
    17<T<18.
    \]
    故选 B.
\end{solution}
\begin{example}[（高考题改编）]
    （单选）数列$a_{n}$各项为正整数且递增，$a_{n+2}=C_{a_{n+1}}^{a_{n}}$，则（~~~~~）

    \begin{tabular}{@{}ll@{}}
        A.$a_n<a_{n-1}+1$&B.$a_1,a_2,a_3$可能成等比数列\\
        C.$a_3a_4<a_5$&D.$a_3,a_4,a_5$可能成等比数列
    \end{tabular}
\end{example}
\begin{solution}
由 $a_n$ 递增可知 $a_n\ge a_{n-1}+1$，故 A 不成立。若 B 或 D 成立，则存在某个 $n$ 使得 $a_n,a_{n+1},a_{n+2}$ 成等比数列，于是
\[
    a_na_{n+2} = a_nC_{a_{n+1}}^{a_{n}}=a_{n+1}C_{a_{n+1}-1}^{a_{n}-1}=a_{n+1}^2\\\Rightarrow a_{n+1}=C_{a_{n+1}-1}^{a_{n}-1}
\]

当$a_n=1$时，代入表达式得到$a_{n+1}=C_{a_{n+1}-1}^{0}=1=a_n$，与数列递增矛盾；

当$a_n=2$时，代入表达式得到$a_{n+1}=C_{a_{n+1}-1}^{1}=a_{n+1}-1<a_{n+1}$，矛盾；

当 $a_n>2$ 时，令 $N=a_{n+1}-1,\ r=a_n-1$，则 $2\le r\le N-1$，且需要
\[
C_N^r=N+1.
\]
若 $r=N-1$，则 $C_N^r=N<N+1$；若 $2\le r\le N-2$，则
\[
C_N^r\ge C_N^2=\frac{N(N-1)}2>N+1.
\]
所以 B、D 均不成立。

验证 C 项。若 $a_1=1$，则 $a_3=C_{a_2}^{1}=a_2$，与数列递增矛盾，故 $a_1\ne1$。于是 $a_1\ge2,\ a_2\ge3$。若 $a_1=2,\ a_2=3$，则 $a_3=C_3^2=3$，仍不递增；故实际有 $a_2\ge4$，从而
\[
a_3=C_{a_2}^{a_1}\ge C_4^2=6.
\]
若 $a_3=a_2+1$，则 $C_{a_2}^{a_1}=a_2+1$。当 $a_1=a_2-1$ 时左端为 $a_2$；当 $2\le a_1\le a_2-2$ 时左端至少为 $C_{a_2}^{2}>a_2+1$，均不可能。因此 $a_3\ge a_2+2$。
此时 $4\le a_2\le a_3-2$。由二项式系数的对称性与前半段递增性，
\[
a_4=C_{a_3}^{a_2}\ge C_{a_3}^{2}>2a_3+1.
\]
又
\[
a_5=C_{a_4}^{a_3}=\frac{a_4}{a_3}C_{a_4-1}^{a_3-1}.
\]
由 $a_4>2a_3+1$ 得 $a_4-1\ge2a_3+1$。固定下标 $a_3-1$ 时，上标增大组合数不减；又 $2\le a_3-1\le \lfloor(2a_3+1)/2\rfloor$，所以在前半段可与第二项比较，故
\[
C_{a_4-1}^{a_3-1}\ge C_{2a_3+1}^{a_3-1}\ge C_{2a_3+1}^{2}=a_3(2a_3+1)>a_3^2.
\]
代回即得
\[
a_5>a_3a_4.
\]
因此选 C。
\end{solution}
\begin{example}
    已知$\triangle{ABC}$中，$A=3B=9C$，则$\cos A\cos B+\cos B\cos C+\cos C\cos A=$\underline{\hspace{1cm}}.
\end{example}
\begin{solution}
    解得$A=\dfrac{9\pi}{13},B=\dfrac{3\pi}{13},C=\dfrac{\pi}{13}$，用积化和差整理：
\begin{align*}
&\cos A\cos B+\cos B\cos C+\cos C\cos A\\
&=\dfrac12(\cos(A+B)+\cos(A-B)+\cos(B+C)+\cos(B-C)+\cos(A+C)+\cos(A-C))\\
&=\dfrac12(\cos\dfrac{12\pi}{13}+\cos\dfrac{6\pi}{13}+\cos\dfrac{4\pi}{13}+\cos\dfrac{2\pi}{13}+\cos\dfrac{10\pi}{13}+\cos\dfrac{8\pi}{13})\\
&=\dfrac{1}{2\sin\dfrac{\pi}{13}}\sin\dfrac{\pi}{13}(\cos\dfrac{2\pi}{13}+\cos\dfrac{4\pi}{13}+\cos\dfrac{6\pi}{13}+\cos\dfrac{8\pi}{13}+\cos\dfrac{10\pi}{13}+\cos\dfrac{12\pi}{13})\\
\intertext{为把余弦和化成望远镜和，乘以 $\sin\dfrac{\pi}{13}$，再用 $2\sin u\cos v=\sin(u+v)-\sin(v-u)$。}
&=\dfrac{1}{4\sin\dfrac{\pi}{13}}(\sin\dfrac{\pi}{13}-\sin\dfrac{3\pi}{13}+\sin\dfrac{3\pi}{13}-\sin\dfrac{5\pi}{13}+\sin\dfrac{5\pi}{13}-\sin\dfrac{7\pi}{13}\\
&~~~~~~~~~~~~~~~~+\sin\dfrac{7\pi}{13}-\sin\dfrac{9\pi}{13}+\sin\dfrac{9\pi}{13}-\sin\dfrac{11\pi}{13}+\sin\dfrac{11\pi}{13}-\sin\dfrac{13\pi}{13})\\
&=-\dfrac14
\end{align*}
\end{solution}
\begin{theorem}[阿贝尔求和]
    设$B_n$是数列$b_n$的前$n$项和，当$n\ge 2$时，有:\vspace{-5pt}\[
    \sum_{i=1}^na_ib_i=a_nB_n-\sum_{i=1}^{n-1}(a_{i+1}-a_i)B_i\]
\end{theorem}
\begin{myproof}
    当$n\ge 2$时，有
    \begin{align*}
        \sum_{i=1}^na_ib_i=&a_1b_1+\sum_{i=2}^{n}a_i(B_i-B_{i-1})\\
        =&a_1b_1+\sum_{i=2}^{n}a_iB_i-\sum_{i=2}^{n}a_iB_{i-1}\\
        =&\sum_{i=1}^{n}a_iB_i-\sum_{i=1}^{n-1}a_{i+1}B_i\\
        =&a_nB_n-\sum_{i=1}^{n-1}(a_{i+1}-a_i)B_i
    \end{align*}
\end{myproof}
\begin{example}[（来自Fiddie）]
    设数列$\{a_n\}$的各项均为实数，且当$n\ge 2$时，$a_{n+1}+a_{n-1}=|a_n|$.证明：

    （1）存在大于$1$的正整数$m$使得$a_m\le 0$

    （2）存在正整数$m$使得$a_m\le 0,~a_{m+1}\le 0$

    （3）$a_n=a_{n+9}$ 
\end{example}
\begin{solution}
    (1)当$n\ge 2$时，由\[
    \begin{cases}a_{n+1}+a_{n-1}=|a_n|\\a_{n+2}+a_{n}=|a_{n+1}|\end{cases}\Rightarrow a_{n+2}+a_{n+1}+a_n+a_{n-1}=|a_n|+|a_{n+1}|
    \]\[
    \Rightarrow a_{n+2}+a_{n-1}=|a_n|-a_n+|a_{n+1}|-a_{n+1}\ge 0
    \]
    若所有 $a_n(n>1)$ 都为正，取 $n\ge3$，则上式化为 $a_{n+2}+a_{n-1}=0$。但此时 $a_{n+2}>0,\ a_{n-1}>0$，矛盾，故存在大于$1$的正整数$m$使得$a_m\le 0$
 
    \noindent(2)已证存在大于$1$的整数$m$使得$a_m\le 0$。现假设不存在正整数$k$使得$a_k\le 0,~a_{k+1}\le 0$。取满足 $m>1,\ a_m\le0$ 的最小正整数 $m$，则
    \[
    a_1,a_2,\ldots,a_{m-1}>0,\quad a_m\le 0,\quad a_{m+1}>0.
    \]
    由递推式逐项计算：\vspace{-10pt}
    \begin{align*}
    a_{m+1}+a_{m-1}=|a_m|&=-a_m\Rightarrow a_{m+1}=-a_m-a_{m-1}\\
    a_{m+2}+a_m=|a_{m+1}|&=a_{m+1}\Rightarrow a_{m+2}=a_{m+1}-a_m>0+0=0\\
    a_{m+3}+a_{m+1}=|a_{m+2}|&=a_{m+2}\Rightarrow a_{m+3}=a_{m+2}-a_{m+1}=a_{m+1}-a_m-a_{m+1}=-a_m\geq 0\\
    a_{m+4}+a_{m+2}=|a_{m+3}|&=-a_{m}\Rightarrow a_{m+4}=-a_{m}-a_{m+2}=-a_{m+1}<0\\
    a_{m+5}+a_{m+3}=|a_{m+4}|&=a_{m+1}\Rightarrow a_{m+5}=a_{m+1}-a_{m+3}=a_m+a_{m+1}=-a_{m-1}<0
    \end{align*}
    这与假设中不存在相邻两项同为非正矛盾，故存在正整数$k$使得$a_k\le 0,~a_{k+1}\le 0$.

\noindent(3)由第二问可取相邻两项同为非正，记
\[
u=a_m\le0,\qquad v=a_{m+1}\le0.
\]
递推式可改写为
\[
a_{n+1}=|a_n|-a_{n-1}.
\]
从 $a_m=u,\ a_{m+1}=v$ 出发逐项推出
\[
\begin{array}{c|ccccccccccc}
\text{下标} & m&m+1&m+2&m+3&m+4&m+5&m+6&m+7&m+8&m+9&m+10\\
\hline
\text{值} & u&v&-u-v&-u-2v&-v&u+v&-u&-2u-v&-u-v&u&v
\end{array}
\]
由于 $u\le0,\ v\le0$，有
\[
-u-v\ge0,\quad -u-2v\ge0,\quad -v\ge0,\quad u+v\le0,\quad -u\ge0,\quad -2u-v\ge0.
\]
这些符号与递推式中绝对值的取法一致，因此表中每一项都由前两项唯一推出，特别地
\[
(a_{m+9},a_{m+10})=(a_m,a_{m+1}).
\]
从这对相邻项向后递推，得到 $a_{n+9}=a_n\ (n\ge m)$；再由
\[
a_{n-1}=|a_n|-a_{n+1}
\]
向前递推，得到 $n<m$ 时同样成立。因此 $a_{n+9}=a_n$。
\end{solution}
\begin{example}[（南通9调14题）]
    已知$x,y$满足$(\sqrt{x^2+1}-x)(\sqrt{y^2+4}-y)=2$，则$4^{x}+2^{y-1}$的最小值为\underline{\hspace{1cm}}.
\end{example}
\begin{solution}
    作变量代换：\[
    \begin{cases}
        m=\sqrt{x^2+1}-x\Rightarrow x=\dfrac{1}{2m}-\dfrac{m}{2}\\[8pt]
        n=\sqrt{y^2+4}-y\Rightarrow y=\dfrac{2}{n}-\dfrac{n}{2}\end{cases}
    \]
    再由$mn=2\Rightarrow n=\dfrac{2}{m}$：
    \[y=\dfrac{2}{n}-\dfrac{n}{2}=m-\dfrac{1}{m}\Rightarrow y=-2x\]
    所以：
    \[
    4^{x}+2^{y-1}=4^{x}+\dfrac{1}{2\times 4^x}\geq 2\sqrt{\dfrac{1}{2}}=\sqrt2.
    \]
    等号成立时 $4^x=\frac1{\sqrt2}$，即 $x=-\frac14$，此时由 $y=-2x$ 得 $y=\frac12$，满足原条件。
\end{solution}
\newpage
\begin{example}[（深圳中学2025届二轮一阶）]
    \noindent$\triangle{ABC}$中，若
    \[\begin{cases}\overrightarrow{AD}=\dfrac{\lambda}{\lambda+1}\overrightarrow{AC}\\
        \overrightarrow{AE}=\dfrac{\mu}{\mu+1}\overrightarrow{AB}\end{cases}\]
    \noindent 则连接$BD,CE$得到交点$Q$，任取$\triangle{ABC}$所在平面内某一点$P$，有
    \[
    PQ^2=\dfrac{PA^2+\mu PB^2+\lambda PC^2}{1+\lambda+\mu}-\dfrac{\lambda\mu BC^2+\lambda AC^2+\mu AB^2}{\left(1+\lambda+\mu\right)^2}
    \]
\end{example}
\begin{solution}
      因为 $Q$ 同时位于 $BD$ 与 $CE$ 上，先分别用这两条线段上的仿射参数表示 $\overrightarrow{AQ}$。设
      \[
      \begin{cases}\overrightarrow{AQ}=x\overrightarrow{AD}+(1-x)\overrightarrow{AB}=x\dfrac{\lambda}{1+\lambda}\overrightarrow{AC}+(1-x)\overrightarrow{AB}\\[10pt]
      \overrightarrow{AQ}=y\overrightarrow{AE}+(1-y)\overrightarrow{AC}=y\dfrac{\mu}{1+\mu}\overrightarrow{AB}+(1-y)\overrightarrow{AC}\end{cases}\Rightarrow \begin{cases}x=\dfrac{\lambda+1}{1+\lambda+\mu}\\[8pt]y=\dfrac{\mu+1}{1+\lambda+\mu}\end{cases}
      \]
      化为形如$x\overrightarrow{QA}+y\overrightarrow{QB}+z\overrightarrow{QC}=0$的方程:
      \begin{align*}
          \Rightarrow \overrightarrow{AQ}&=\dfrac{\lambda}{1+\lambda+\mu}\overrightarrow{AC}+\dfrac{\mu}{1+\lambda+\mu}\overrightarrow{AB}\\
          &=\dfrac{\lambda}{1+\lambda+\mu}(\overrightarrow{AQ}+\overrightarrow{QC})+\dfrac{\mu}{1+\lambda+\mu}(\overrightarrow{AQ}+\overrightarrow{QB})\\
          &=\dfrac{\lambda+\mu}{1+\lambda+\mu}\overrightarrow{AQ}+\dfrac{\lambda}{1+\lambda+\mu}\overrightarrow{QC}+\dfrac{\mu}{1+\lambda+\mu}\overrightarrow{QB}\\
          \Rightarrow \overrightarrow{QA}&+\lambda\overrightarrow{QC}+\mu\overrightarrow{QB}=\vec{0}\\
          \Rightarrow \overrightarrow{PA}&-\overrightarrow{PQ}+\lambda(\overrightarrow{PC}-\overrightarrow{PQ})+\mu(\overrightarrow{PB}-\overrightarrow{PQ})=\vec{0}\\
          \Rightarrow \overrightarrow{PA}&+\lambda\overrightarrow{PC}+\mu\overrightarrow{PB}=(1+\lambda+\mu)\overrightarrow{PQ}
      \end{align*}
    平方得：\[
PA^2+\lambda^2PC^2+\mu^2PB^2+2\lambda\overrightarrow{PA}\cdot\overrightarrow{PC}+2\mu\overrightarrow{PA}\cdot\overrightarrow{PB}+2\lambda\mu\overrightarrow{PC}\cdot\overrightarrow{PB}=(1+\lambda+\mu)^2PQ^2\]
由
\[
2\overrightarrow{PA}\cdot\overrightarrow{PC}=PA^2+PC^2-AC^2,\quad
2\overrightarrow{PA}\cdot\overrightarrow{PB}=PA^2+PB^2-AB^2,\quad
2\overrightarrow{PC}\cdot\overrightarrow{PB}=PC^2+PB^2-BC^2
\]
代入得
\[\begin{cases}2\lambda\overrightarrow{PA}\cdot\overrightarrow{PC}=\lambda(PA^2+PC^2-AC^2)\\
    2\mu\overrightarrow{PA}\cdot\overrightarrow{PB}=\mu(PA^2+PB^2-AB^2)\\
    2\lambda\mu\overrightarrow{PC}\cdot\overrightarrow{PB}=\lambda\mu(PC^2+PB^2-BC^2)
\end{cases}\]
    整理得：\[
    PQ^2=\dfrac{PA^2+\mu PB^2+\lambda PC^2}{1+\lambda+\mu}-\dfrac{\lambda\mu BC^2+\lambda AC^2+\mu AB^2}{\left(1+\lambda+\mu\right)^2}\]
\end{solution}
\begin{conclusion}[平方关系]
    \noindent 平面内给定$\triangle{ABC}$，若点$Q$满足\vspace{-10pt}
    \[\overrightarrow{QA}+\lambda\overrightarrow{QC}+\mu\overrightarrow{QB}=\vec{0}\vspace{-10pt}\]则任取$\triangle{ABC}$所在平面内某一点$P$，有\[PQ^2=\dfrac{PA^2+\mu PB^2+\lambda PC^2}{1+\lambda+\mu}-\dfrac{\lambda\mu BC^2+\lambda AC^2+\mu AB^2}{\left(1+\lambda+\mu\right)^2}\]
\end{conclusion}
\begin{example}[奔驰定理]
    \noindent 已知平面直角坐标系$xOy$中有一个$\triangle{ABC}$，以及平面内任意一点$P$，则有：\[
    \begin{vmatrix}
    x_B & y_B & 1\\
    x_C & y_C & 1\\
    x_P & y_P & 1\\
\end{vmatrix}\overrightarrow{PA}+\begin{vmatrix}
    x_C & y_C & 1\\
    x_A & y_A & 1\\
    x_P & y_P & 1\\
\end{vmatrix}\overrightarrow{PB}+\begin{vmatrix}
    x_A & y_A & 1\\
    x_B & y_B & 1\\
    x_P & y_P & 1\\
\end{vmatrix}\overrightarrow{PC}=\vec{0}\]
这也可以写成\[S_{\triangle{PBC}}\overrightarrow{PA}+S_{\triangle{PAC}}\overrightarrow{PB}+S_{\triangle{PAB}}\overrightarrow{PC}=\vec{0}\]
这里的三角形面积取有向面积。计算时按字母顺序确定方向，与另外两个方向相反的面积取负值。
\end{example}
\begin{solution}
记有向面积
\[
[XYZ]=\frac12\begin{vmatrix}
x_X & y_X & 1\\
x_Y & y_Y & 1\\
x_Z & y_Z & 1
\end{vmatrix}.
\]
点 $P$ 关于 $\triangle ABC$ 的有向面积坐标满足
\[
[PBC]+[APC]+[ABP]=[ABC],
\]
且
\[
[PBC]A+[APC]B+[ABP]C=[ABC]P.
\]
两边同减去 $[ABC]P$，便得
\[
[PBC]\overrightarrow{PA}+[APC]\overrightarrow{PB}+[ABP]\overrightarrow{PC}=\vec0.
\]
把有向面积写成行列式形式，即为题中公式。
\end{solution}
\begin{example}[外心向量关系]
    已知平面直角坐标系$xOy$中有一个$\triangle{ABC}$，若其外心为$U(x,y)$，则$U$满足
    \[
    2U\cdot(B-A)=|B|^2-|A|^2,\qquad
    2U\cdot(C-A)=|C|^2-|A|^2.
    \]
\end{example}
\begin{solution}
外心到三顶点距离相等，因此
\[
|U-A|^2=|U-B|^2,\qquad |U-A|^2=|U-C|^2.
\]
展开第一式：
\[
|U|^2-2U\cdot A+|A|^2=|U|^2-2U\cdot B+|B|^2,
\]
移项即得
\[
2U\cdot(B-A)=|B|^2-|A|^2.
\]
第二式同理。
\end{solution}
\begin{example}[三角形外心坐标]
    已知平面直角坐标系$xOy$中有一个$\triangle{ABC}$，则其外心的坐标为\[
    \left(\dfrac{\begin{vmatrix}
    OA^2 & y_A & 1\\
    OB^2 & y_B & 1\\
    OC^2 & y_C & 1\\
\end{vmatrix}}{2\begin{vmatrix}
    x_A & y_A & 1\\
    x_B & y_B & 1\\
    x_C & y_C & 1\\
\end{vmatrix}},\dfrac{\begin{vmatrix}
    x_A & OA^2 & 1\\
    x_B & OB^2 & 1\\
    x_C & OC^2 & 1\\
\end{vmatrix}}{2\begin{vmatrix}
    x_A & y_A & 1\\
    x_B & y_B & 1\\
    x_C & y_C & 1\\
\end{vmatrix}}\right)
\]\end{example}
\begin{solution}
设外心为 $U(x,y)$。由上一题，
\[
\begin{cases}
2x(x_B-x_A)+2y(y_B-y_A)=OB^2-OA^2,\\
2x(x_C-x_A)+2y(y_C-y_A)=OC^2-OA^2.
\end{cases}
\]
把它看作关于 $x,y$ 的二元一次方程组。记
\[
\Delta=\begin{vmatrix}
x_B-x_A&y_B-y_A\\
x_C-x_A&y_C-y_A
\end{vmatrix}
=\begin{vmatrix}
x_A&y_A&1\\
x_B&y_B&1\\
x_C&y_C&1
\end{vmatrix}.
\]
由克莱姆法则，
\[
x=\frac{(OB^2-OA^2)(y_C-y_A)-(OC^2-OA^2)(y_B-y_A)}{2\Delta}.
\]
分子等于把
\[
\begin{vmatrix}
OA^2&y_A&1\\
OB^2&y_B&1\\
OC^2&y_C&1
\end{vmatrix}
\]
的第二、三行分别减去第一行后得到的二阶行列式。$y$ 坐标同理：
\[
y=\frac{(x_B-x_A)(OC^2-OA^2)-(x_C-x_A)(OB^2-OA^2)}{2\Delta},
\]
对应题中第二个三阶行列式。
\end{solution}
\begin{example}[三角形垂心坐标]
    已知平面直角坐标系$xOy$中有一个$\triangle{ABC}$，则其垂心的坐标为\[
    \left(\dfrac{\begin{vmatrix}
    x_Bx_C+y_By_C & y_A & 1\\
    x_Ax_C+y_Ay_C & y_B & 1\\
    x_Ax_B+y_Ay_B & y_C & 1\\
    \end{vmatrix}}{-\begin{vmatrix}
    x_A & y_A & 1\\
    x_B & y_B & 1\\
    x_C & y_C & 1\\
    \end{vmatrix}},\dfrac{\begin{vmatrix}
    x_A & x_Bx_C+y_By_C & 1\\
    x_B & x_Ax_C+y_Ay_C & 1\\
    x_C & x_Ax_B+y_Ay_B & 1\\
    \end{vmatrix}}{-\begin{vmatrix}
    x_A & y_A & 1\\
    x_B & y_B & 1\\
    x_C & y_C & 1\\
    \end{vmatrix}}\right)\]
\end{example}
\begin{solution}
设垂心为 $H(x,y)$。由 $AH\perp BC$ 与 $BH\perp AC$，
\[
(H-A)\cdot(B-C)=0,\qquad (H-B)\cdot(A-C)=0.
\]
也就是
\[
\begin{cases}
x(x_B-x_C)+y(y_B-y_C)=x_A(x_B-x_C)+y_A(y_B-y_C),\\
x(x_A-x_C)+y(y_A-y_C)=x_B(x_A-x_C)+y_B(y_A-y_C).
\end{cases}
\]
对这个二元一次方程组使用克莱姆法则。其系数行列式为
\[
\begin{vmatrix}
x_B-x_C&y_B-y_C\\
x_A-x_C&y_A-y_C
\end{vmatrix}
=-\begin{vmatrix}
x_A&y_A&1\\
x_B&y_B&1\\
x_C&y_C&1
\end{vmatrix}.
\]
以 $x$ 坐标为例，克莱姆法则的分子为
\[
\begin{vmatrix}
x_A(x_B-x_C)+y_A(y_B-y_C)&y_B-y_C\\
x_B(x_A-x_C)+y_B(y_A-y_C)&y_A-y_C
\end{vmatrix}.
\]
把第一列按点积展开，并作行列式的初等行变换，可化为
\[
\begin{vmatrix}
x_Bx_C+y_By_C&y_A&1\\
x_Ax_C+y_Ay_C&y_B&1\\
x_Ax_B+y_Ay_B&y_C&1
\end{vmatrix}.
\]
$y$ 坐标的化简对应进行，只是把上述行列式中的第二列换成点积列，得到题中公式。
\end{solution}
\begin{example}[三角形的内心坐标]
    已知平面直角坐标系$xOy$中有一个$\triangle{ABC}$，则其内心的坐标为\[
    \left(\dfrac{ax_A+bx_B+cx_C}{a+b+c},\dfrac{ay_A+by_B+cy_C}{a+b+c}\right)
    \]
\end{example}
\begin{solution}
设内心为 $I$，内切圆半径为 $r$。若
\[
a=|BC|,\quad b=|CA|,\quad c=|AB|,
\]
则
\[
[IBC]=\frac12ar,\qquad [ICA]=\frac12br,\qquad [IAB]=\frac12cr.
\]
因此 $I$ 关于 $\triangle ABC$ 的面积坐标为 $a:b:c$，即
\[
I=\frac{aA+bB+cC}{a+b+c}.
\]
写成坐标形式就是题中公式。
\end{solution}
\begin{example}[容斥原理练习]
    某学校举办比赛，有$20$个参赛名额，现在分给$4$个不同的班，保证至少有一个班的名额为$4$个，且每一个班都有名额，则共有\underline{\hspace{1cm}}种分法。
\end{example}
\begin{solution}
    设四个班的名额为$x_1,x_2,x_3,x_4\in N_+$，则分法数就是集合$A_i=\{(x_1,x_2,x_3,x_4)|x_i=4,x_1+x_2+x_3+x_4=20\}$的元素个数，又因为
    \begin{align*}|A_{1}|&=\#\{(x_{2},x_{3},x_{4})\mid x_{2}+x_{3}+x_{4}=16,x_{2},x_{3},x_{4}>0\}=C_{15}^{2}\\|A_{1}\cap A_{2}|&=\#\{(x_{3},x_{4})\mid x_{3}+x_{4}=12,x_{3},x_{4}>0\}=C_{11}^{1}\\|A_{1}\cap A_{2}\cap A_{3}|&=1\\
    |A_1\cup A_2\cup A_3\cup A_4|&=C_4^1|A_1|-C_4^2|A_1\cap A_2|+C_4^3|A_1\cap A_2\cap A_3|\\
        &=C_4^1C_{15}^2-C_4^2C_{11}^1+C_4^3=358\end{align*}
\end{solution}
\newpage
\begin{example}[求和]
    计算$\displaystyle \sum_{i=0}^{n-1}(-1)^{i}C_{n-1}^{i}(i+1)^k$，其中$k<n-1,k\in N_+$
\end{example}
\begin{solution}
    为把交错组合和转化为导数值，取生成函数
    \begin{align*}f(x)=&\displaystyle\sum_{i=0}^{n-1}(-1)^{i}C_{n-1}^{i}x^{i+1}
        =\displaystyle x\sum_{i=0}^{n-1}(-x)^{i}C_{n-1}^{i}\\
        =&\displaystyle x(1-x)^{n-1}
        =(x-1+1)(1-x)^{n-1}
        =(1-x)^{n-1}-(1-x)^n\\
    \Rightarrow f(1)=&\displaystyle\sum_{i=0}^{n-1}(-1)^{i}C_{n-1}^{i}=0\\
     f^{(k)}(x)=&\displaystyle\sum_{i=0}^{n-1}(-1)^{i}C_{n-1}^{i}A_{i+1}^{k}x^{i+1-k}\Rightarrow
     f^{(k)}(1)=\displaystyle\sum_{i=0}^{n-1}(-1)^{i}C_{n-1}^{i}A_{i+1}^{k}
    \end{align*}
    将$A_{i+1}^{k}$和${(i+1)}^{k}$联系起来。考虑把$k$个有标号的球放入$i+1$个有标号的盒子中，允许空盒，则放法数为${(i+1)}^{k}$。若按非空盒子的个数分类，则
    \[\sum_{r=0}^kS(k,r)r!C_{i+1}^r=\sum_{r=0}^kS(k,r)A_{i+1}^r\]
    其中$S(k,r)$表示把$k$个有标号的球放入$r$个无标号非空盒的方法数，$C_{i+1}^r$表示从$i+1$个盒子中选出$r$个，再排列后得到有标号的$r$个盒子。因此
    \[(i+1)^{k}=\sum_{r=0}^kS(k,r)A_{i+1}^r
    \]
    代回原式，得到\begin{align*}
    \sum_{i=0}^{n-1}(-1)^{i}C_{n-1}^{i}(i+1)^k=&\sum_{i=0}^{n-1}\bigg((-1)^{i}C_{n-1}^{i}\bigg(\sum_{r=0}^kS(k,r)A_{i+1}^r\bigg)\bigg)\\=&\sum_{i=0}^{n-1}\sum_{r=0}^kS(k,r)A_{i+1}^r(-1)^{i}C_{n-1}^{i}\\=&\sum_{r=0}^k\sum_{i=0}^{n-1}S(k,r)A_{i+1}^r(-1)^{i}C_{n-1}^{i}=\sum_{r=0}^kS(k,r)f^{(r)}(1)
    \end{align*}
    由$(1-x)^{n-1}-(1-x)^n$的导数可知，$f(x)$只有第$n-1$和$n$阶导数在$x=1$处的值不是$0$，即
    \[\sum_{i=0}^{n-1}(-1)^{i}C_{n-1}^{i}(i+1)^k=0
    \]
\end{solution}
\begin{example}[证明]
    \vspace{-15pt}\begin{align*}
    (1)\quad&\frac{\sin B - \cos A}{\sin A + \cos B} - \frac{\sin A + \cos B}{\sin B - \cos A} = 2\tan(A+B)\\
    (2)\quad&\frac{\sin B - \cos A}{\sin A - \cos B} - \frac{\sin A - \cos B}{\sin B - \cos A} = \frac{\cos B + \sin A}{\cos A + \sin B} - \frac{\cos A + \sin B}{\cos B + \sin A} = 2\tan(A-B)\\
    (3)\quad&\frac{\cos A - \cos B}{\sin A - \sin B} - \frac{\sin A - \sin B}{\cos A - \cos B} = 2\cot(A+B)\\
    (4)\quad&\frac{\cos A - \cos B}{\sin A + \sin B} - \frac{\sin A + \sin B}{\cos A - \cos B} = \frac{\cos A + \cos B}{\sin A - \sin B} - \frac{\sin A - \sin B}{\cos A + \cos B} = 2\cot(A-B)\end{align*}
\end{example}
\begin{solution}
第二式通分整理得：
\begin{align*}
&\dfrac{\sin B - \cos A}{\sin A - \cos B} - \dfrac{\sin A - \cos B}{\sin B - \cos A}\\[1.2ex]
&=\dfrac{(\sin B - \cos A)^2 - (\sin A - \cos B)^2}{(\sin B - \cos A)\left(\sin A - \cos B\right)}\\[1.2ex]
&=\dfrac{\sin^2 B - 2\sin B \cos A + \cos^2 A - \sin^2 A + 2\sin A \cos B - \cos^2 B}{\sin B \sin A - \sin A \cos A - \sin B \cos B + \cos B \cos A}\\[1.2ex]
&=\dfrac{\cos 2A - \cos 2B + 2\left(\sin A \cos B - \sin B \cos A\right)}{\cos\left(A - B\right) - \dfrac{1}{2}\left(\sin 2A + \sin 2B\right)}\\[1.2ex]
&=\dfrac{-2\sin\left(A + B\right)\sin\left(A - B\right) + 2\sin\left(A - B\right)}{\cos\left(A - B\right) - \sin\left(A + B\right)\cos\left(A - B\right)}
=\dfrac{2\left[1 - \sin\left(A + B\right)\right]\sin\left(A - B\right)}{\left[1 - \sin\left(A + B\right)\right]\cos\left(A - B\right)}
\\[1.2ex]&=\dfrac{2\sin\left(A - B\right)}{\cos\left(A - B\right)} = 2\tan\left(A - B\right).
\end{align*}
第三式同样通分可得：
\begin{align*}
&\dfrac{\cos A - \cos B}{\sin A - \sin B} - \dfrac{\sin A - \sin B}{\cos A - \cos B}\\[1.2ex]
&=\dfrac{(\cos A - \cos B)^2 - (\sin A - \sin B)^2}{(\cos A - \cos B)(\sin A - \sin B)}\\[1.2ex]
&=\dfrac{\cos^2 A - 2\cos A\cos B + \cos^2 B - \sin^2 A + 2\sin A\sin B - \sin^2 B}{\sin A\cos A - \sin A\cos B - \sin B\cos A + \sin B\cos B}\\[1.2ex]
&=\dfrac{(\cos^2 A - \sin^2 A) + (\cos^2 B - \sin^2 B) - 2\cos A\cos B + 2\sin A\sin B}{\sin A\cos A + \sin B\cos B - (\sin A\cos B + \sin B\cos A)}\\[1.2ex]
&=\dfrac{\cos 2A + \cos 2B - 2\cos(A + B)}{\dfrac{1}{2}(\sin 2A + \sin 2B) - \sin(A + B)}
=\dfrac{2\cos(A + B)\cos(A - B) - 2\cos(A + B)}{\sin(A + B)\cos(A - B) - \sin(A + B)}\\[1.2ex]
&=\dfrac{2\cos(A + B)[\cos(A - B) - 1]}{\sin(A + B)[\cos(A - B) - 1]}
=\dfrac{2\cos(A + B)}{\sin(A + B)} = 2\cot(A + B).
\end{align*}
第一式可由第二式把 $B$ 换成 $\pi-B$ 得到；第四式中的两种写法，分别由第三式把 $B$ 换成 $-B$ 与 $\pi-B$ 得到。
\end{solution}
\newpage
\begin{example}[（椭圆弦与平行条件）]
    平面直角坐标系$xOy$中，椭圆$\frac{x^2}{16}+\frac{y^2}{a^2}=1,(a>4)$，椭圆的弦$AB$过点$D(-2,0)$，连接$OA,OB$，线段$OB$交圆$(x+1)^2+y^2=4$于$C$，若$DC$平行于$OA$，求$a=$\underline{\hspace{1.2cm}}.
\end{example}
\begin{solution}
    设$B(x_1,y_1),A(x_2,y_2),C(x_3,y_3)$。由于 $A,D,B$ 共线且 $D(-2,0)$，可用有向比记录 $D$ 分割 $AB$ 的位置，记
    \[
    -\lambda=\dfrac{x_1+2}{x_2+2}=\dfrac{y_1}{y_2}.
    \]
    则
    \[\begin{cases}\dfrac{x_1^2}{16}+\dfrac{y_1^2}{a^2}=1&\Rightarrow \dfrac{x_1^2-\lambda^2x_2^2}{16}+\dfrac{y_1^2-\lambda^2y_2^2}{a^2}=1-\lambda^2\\[1.2ex]\dfrac{\lambda^2x_2^2}{16}+\dfrac{\lambda^2y_2^2}{a^2}=\lambda^2&\Rightarrow \dfrac{(x_1+\lambda x_2)(x_1-\lambda x_2)}{16(1+\lambda)(1-\lambda)}+\dfrac{(y_1+\lambda y_2)(y_1-\lambda y_2)}{a^2(1+\lambda)(1-\lambda)}=1\end{cases}\]
    代入$x_1+\lambda x_2=(-2)(1+\lambda),y_1+\lambda y_2=0$得
    \begin{align*}\dfrac{-(x_1-\lambda x_2)}{8(1-\lambda)}=1&\Rightarrow \lambda=\dfrac{x_1+8}{x_2+8}\Rightarrow -\lambda=\dfrac{x_1+2}{x_2+2}=\dfrac{y_1}{y_2}=\dfrac{x_1+8}{-x_2-8}\\
    &\Rightarrow y_2=\dfrac{-3y_1}{x_1+5},\quad x_2=\dfrac{-3(x_1+2)}{x_1+5}-2=\dfrac{-5x_1-16}{x_1+5}
\end{align*}
设$y_1=kx_1$，则代入$\dfrac{x^2}{16}+\dfrac{y^2}{a^2}=1$得
\[
\dfrac{x_1^2}{16}+\dfrac{k^2x_1^2}{a^2}=1\Rightarrow \begin{cases}x_1=\pm\dfrac{4a}{\sqrt{a^2+16k^2}}\\[1.2ex]y_1=\pm\dfrac{4ak}{\sqrt{a^2+16k^2}}\end{cases}
\]
再代入圆方程，得到
\[
(x+1)^2+k^2x^2=4\Rightarrow (k^2+1)x^2+2x-3=0\Rightarrow \begin{cases}x_3=\dfrac{-1\pm\sqrt{3k^2+4}}{k^2+1}\\y_3=\dfrac{-k\pm k\sqrt{3k^2+4}}{k^2+1}\end{cases}.
\]
点 $B,C$ 都在直线 $OB$ 上，所取正负号需对应同一射线方向，即上下的正负号同时取正或同时取负：
\begin{align*}
    k_{AO}&=\dfrac{3y_1}{5x_1+16}=\dfrac{\dfrac{\pm 12ak}{\sqrt{a^2+16k^2}}}{\pm\dfrac{20a}{\sqrt{a^2+16k^2}}+16}=\dfrac{\pm 3ak}{4\sqrt{a^2+16k^2}\pm 5a}\\[1.2ex]
    k_{DC}&=\dfrac{\dfrac{-k\pm k\sqrt{3k^2+4}}{k^2+1}}{\dfrac{-1\pm\sqrt{3k^2+4}}{k^2+1}+2}=\dfrac{-k\pm k\sqrt{3k^2+4}}{2k^2+1\pm \sqrt{3k^2+4}}
\end{align*}
记 $r=\sqrt{3k^2+4}$。因为
\[
2k^2+1\pm r=\frac{(r\mp1)(2r\pm5)}3,
\]
所以
\[
\frac{-1\pm r}{2k^2+1\pm r}=\frac{\pm3}{2r\pm5}.
\]
由 $k_{AO}=k_{CD}$ 得
\[
\frac{\pm3}{2r\pm5}=\frac{\pm3a}{4\sqrt{a^2+16k^2}\pm5a}.
\]
交叉相乘并消去相同项，得到
\[
\frac{\sqrt{a^2+16k^2}}{\sqrt{3k^2+4}}=\frac a2.
\]
题目要求由过 $D$ 的动弦得到同一个参数 $a$，所以平行条件不能依赖某一个特定斜率。于是上式应对这族弦对应的可取 $k$ 都成立；把左端视为 $k^2$ 的函数，只能是常数，故
\[
\frac{a^2+16k^2}{3k^2+4}=\frac{a^2}{4}
\]
与 $k$ 无关。比较 $k^2$ 项，得
\[
\frac{16}{3}=\frac{a^2}{4},
\]
从而 $a=\dfrac{8}{\sqrt{3}}>4$。此时代回上式有
\[
\sqrt{a^2+16k^2}=\frac a2\sqrt{3k^2+4},
\]
故斜率等式成立，所以
\[
a=\dfrac{8}{\sqrt{3}}.
\]
\end{solution}
\begin{example}[（杭二模拟）]
    \noindent 已知椭圆$\dfrac{x^2}{4}+\dfrac{y^2}{3}=1$，设其左焦点为$F(-1,0)$\vspace{10pt}\\
    （1）若椭圆上的两点$M,N$满足直线$MF,NF$关于$x$轴对称，且直线$MN$斜率存在，证明直线$MN$过定点.\\
    （2）若过$B(-2,2)$的一条直线交椭圆于$M,N$，$A(2,0)$，连接直线$AM,AN$交直线$OB$于$P,Q$，求$\dfrac{|OP|}{|OQ|}$的值.
\end{example}
\begin{solution}
    （1）设过焦点 $F(-1,0)$ 且关于 $x$ 轴对称的两条直线为 $x+1+ty=0$ 与 $x+1-ty=0$。又设直线$MN$的方程为$x=my+n$，则
    \[\begin{cases}(x+1+ty)(x+1-ty)=0\\x=my+n\end{cases}\]
与
\[
\begin{cases}\dfrac{x^2}{4}+\dfrac{y^2}{3}=1\\x=my+n\end{cases}
\]
对应的 $y$ 根同为 $M,N$ 的纵坐标，故两个关于 $y$ 的二次方程成比例。分别化为
\[
(m^2-t^2)y^2+2m(n+1)y+(n+1)^2=0,
\]
\[
(3m^2+4)y^2+6mny+3n^2-12=0.
\]
比较一次项与常数项，得
    \[\begin{cases} 2m(n+1)(3m^2 + 4) = 6mn(m^2 - t^2)\\(n+1)^2(3m^2 + 4) = (3n^2-12)(m^2 - t^2)\end{cases}\]
由于直线 $MN$ 的斜率存在，$m\ne0$；若 $n=0$，第一式左端非零而右端为零，不合，故 $n\ne0$。由第一式解得
\[
m^2-t^2=\frac{(n+1)(3m^2+4)}{3n}.
\]
代入第二式并约去非零因子 $3m^2+4$，得
\[
(n+1)^2=\frac{(3n^2-12)(n+1)}{3n}.
\]
非退化情形下 $n+1\ne0$，否则前一组退化二次式会恒为零，不能给出两点 $M,N$。于是
\[
n(n+1)=n^2-4,
\]
解得 $n=-4$。因此直线$MN$过定点$(-4,0)$。若直接比较$\dfrac{1}{x_1}+\dfrac{1}{x_2}$，也得到同样的结论。
    （2）设$M\bigg(\dfrac{2(1-t_1^2)}{1+t_1^2},\dfrac{2\sqrt{3}t_1}{1+t_1^2}\bigg),N\bigg(\dfrac{2(1-t_2^2)}{1+t_2^2},\dfrac{2\sqrt{3}t_2}{1+t_2^2}\bigg)$，则\[MA:x=\dfrac{x_1-2}{y_1}y+2=\dfrac{\frac{2(1-t_1^2)}{1+t_1^2}-2}{\frac{2\sqrt{3}t_1}{1+t_1^2}}y+2=-\dfrac{2t_1}{\sqrt3}y+2~~~~~~NA:x=-\dfrac{2t_2}{\sqrt3}y+2\]由直线$MA,NA$与$y=-x$联立得到
    \[y_P=\dfrac{2\sqrt3}{2t_1-\sqrt3},y_Q=\dfrac{2\sqrt3}{2t_2-\sqrt3}\Rightarrow \dfrac{|OP|}{|OQ|}=\left|\dfrac{2t_2-\sqrt3}{2t_1-\sqrt3}\right|\]
    由题设，直线 $MN$ 经过 $B(-2,2)$。椭圆参数式下，过参数 $t_1,t_2$ 两点的直线可写为
    \[
    \frac{x}{2}\frac{1-t_1t_2}{1+t_1t_2}+\frac{y}{\sqrt3}\frac{t_1+t_2}{1+t_1t_2}=1.
    \]
    把 $B$ 代入，得到
    \[\dfrac{-2}{2}\dfrac{1-t_1t_2}{1+t_1t_2}+\dfrac{2}{\sqrt3}\dfrac{t_1+t_2}{1+t_1t_2}=1\Rightarrow t_1+t_2=\sqrt3\]
    因而 $2t_2-\sqrt3=-(2t_1-\sqrt3)$，得到$\dfrac{|OP|}{|OQ|}=1$.
\end{solution}

\begin{conclusion}[比值代换法（坐标式）]
\noindent 椭圆$\dfrac{x^2}{a^2}+\dfrac{y^2}{b^2}=1$上两点$A(x_1,y_1),B(x_2,y_2)$与不在椭圆上的点$P(x_0,y_0)$满足$\overrightarrow{AP}=\lambda\overrightarrow{PB}$，则
\[\begin{cases}\lambda=\dfrac{x_1-x_0}{x_0-x_2}=\dfrac{y_1-y_0}{y_0-y_2}=-\dfrac{\dfrac{x_0x_1}{a^2}+\dfrac{y_0y_1}{b^2}-1}{\dfrac{x_0x_2}{a^2}+\dfrac{y_0y_2}{b^2}-1}\\(1+\lambda)^2\bigg(\dfrac{x_0^2}{a^2}+\dfrac{y_0^2}{b^2}-1\bigg)=2\lambda\bigg(\dfrac{x_1x_2}{a^2}+\dfrac{y_1y_2}{b^2}-1\bigg)\end{cases}\]
\end{conclusion}
\begin{myproof}
记
\[
T(M,N)=\frac{x_Mx_N}{a^2}+\frac{y_My_N}{b^2}-1.
\]
由
\[
P=\frac{A+\lambda B}{1+\lambda}
\]
可得
\[
T(P,A)=\frac{T(A,A)+\lambda T(B,A)}{1+\lambda}
=\frac{\lambda T(A,B)}{1+\lambda},
\]
\[
T(P,B)=\frac{T(A,B)+\lambda T(B,B)}{1+\lambda}
=\frac{T(A,B)}{1+\lambda}.
\]
于是
\[
\lambda=\frac{T(P,A)}{T(P,B)}.
\]
再由 $P=\frac{A+\lambda B}{1+\lambda}$ 代入 $T(P,P)$，并利用 $T(A,A)=T(B,B)=0$，即可推出
\[
T(P,P)T(A,B)=2T(P,A)T(P,B).
\]
\end{myproof}

\begin{conclusion}[比值代换法（极化式）]
\noindent 记
\[
T(M, N) = \dfrac{x_M x_N}{a^2} + \dfrac{y_M y_N}{b^2} - 1.
\]
若椭圆上两点$A(x_1,y_1), B(x_2,y_2)$与不在椭圆上的点$P(x_0,y_0)$共线，且满足$\overrightarrow{AP} = \lambda \overrightarrow{PB}$，则
\[
\lambda = \dfrac{x_1-x_0}{x_0-x_2} = \dfrac{T(P,A)}{T(P,B)}.
\]
并且
\[
T(A,B)=T(P,A)+T(P,B),\qquad T(P,P)\,T(A,B)=2\,T(P,A)\,T(P,B).
\]
消去$\lambda$后可得
\[
\dfrac{2}{T(P,P)} = \dfrac{1}{T(P,A)} + \dfrac{1}{T(P,B)}.
\]
\end{conclusion}


\begin{example}[比较关系]
    椭圆$\dfrac{x^2}{a^2}+\dfrac{y^2}{b^2}=1$上两点$A(x_1,y_1),B(x_2,y_2)$，比较$(x_1y_2+x_2y_1)\left(\frac{x_1x_2}{a^2}+\frac{y_1y_2}{b^2}\right)$和$x_1y_1+x_2y_2$的大小关系，并说明结论.
\end{example}
\begin{solution}
令
\[
X_i=\frac{x_i}{a},\qquad Y_i=\frac{y_i}{b}\quad(i=1,2).
\]
则 $X_i^2+Y_i^2=1$。设
\[
X_i=\cos\theta_i,\qquad Y_i=\sin\theta_i.
\]
于是
\[
\frac{x_1x_2}{a^2}+\frac{y_1y_2}{b^2}
=\cos(\theta_1-\theta_2),
\]
且
\[
x_1y_2+x_2y_1=ab\sin(\theta_1+\theta_2).
\]
另一方面，
\[
x_1y_1+x_2y_2
=ab(\cos\theta_1\sin\theta_1+\cos\theta_2\sin\theta_2)
=ab\sin(\theta_1+\theta_2)\cos(\theta_1-\theta_2).
\]
两式相同，故
\[
(x_1y_2+x_2y_1)\left(\frac{x_1x_2}{a^2}+\frac{y_1y_2}{b^2}\right)
=x_1y_1+x_2y_2.
\]
\end{solution}

\begin{conclusion}[参数方程]
    过椭圆$\dfrac{x^2}{a^2}+\dfrac{y^2}{b^2}=1$上两点$A\left(a\dfrac{1-t_1^2}{1+t_1^2},b\dfrac{2t_1}{1+t_1^2}\right),B\left(a\dfrac{1-t_2^2}{1+t_2^2},b\dfrac{2t_2}{1+t_2^2}\right)$的直线方程为\[\dfrac{x(1-t_1t_2)}{a(1+t_1t_2)}+\dfrac{y(t_1+t_2)}{b(1+t_1t_2)}=1\]
\end{conclusion}


